\section{Test set}
\label{sec:horizons:test}

The locked period of 24--26 March was evaluated once, after every choice of the preceding
sections had been made, with the protocol below.  The numbers of this section are the final
ones of the thesis; where they differ from the validation day, \cref{ch:conclusions}
follows them.

\paragraph{Protocol.}  The stateful models are the saved checkpoints of the validated runs,
weights only, run over the test stream in evaluation mode with no learning: three warm-up
days, 21--23 March, rebuild the per-market states, the hierarchical memory and the book
buffers that a checkpoint does not hold, and the three days that follow are recorded.  The
evaluation mode was checked on the validation day first: re-run on 13 March after a two-day
warm-up, the checkpoints reproduce the hard predictions of their original validation files
on 95 to 99.99\% of the samples and their balanced accuracy to the third decimal, with the
lowest agreement at six and twelve hours, where the state lives longest, and a seven-day
warm-up moves the agreement by less than one point; the residual disagreement comes from the
original validation states, which were built while the weights were still changing, not
from the length of the warm-up.  The same re-runs supply the validation probabilities of
the older runs, so that at every horizon the comparison is between calibrated models.  Every
model receives two logit offsets (on \textsc{down} and \textsc{up}) chosen once on the whole
validation day, on the grid of $\pm 2$ in steps of $0.1$, and applied unchanged to the test
days: unlike the cross-fitted offsets of \cref{sec:models:protocol}, this is a calibration
fixed on one day and used eleven days later, and how well it transfers is itself one of the
results.  The stateless control is LightGBM refitted exactly as in validation (fit rows of
7--11 March, early stopping on 12 March, seed 0) and evaluated on the samples of the stateful
test run of the same horizon, so that every comparison is paired; the tabular networks are
trained on the same table; Mean Reversion is computed on the same samples; the intervals are
95\% hour-block bootstraps over the test hours.  The samples are the labelled samples of the
test stream at each horizon, between 12.5 and 17.1 million (60 to 72 hours; the last twelve
hours of 26 March have no label at the longest horizon).  The zero-shot forecaster is
evaluated as in \cref{sec:external:tsfm}, on 60{,}000 random samples of the test run of each
horizon.  Last, a second seed of each of the two one-minute stateful models, alone and as
the residual of the trees, trained after this evaluation with the recipe of the validated
runs unchanged and evaluated the same way, measures the seed noise of the one-minute
results; nothing in them was chosen on the test days.

\paragraph{Coverage of the test days.}  The three days are not three ordinary days.  The
collector holds 14.1, 12.5 and 6.6 million order-book snapshots for 24, 25 and 26 March,
against 6.6 to 19.7 million per day over the whole window and 12.0 million on the
validation day: every hour of 26 March is present, but each holds about 40\% of the usual
rows.  Whether the gap is the collector's or the exchange's could not be established; the
results are therefore given per day as well as pooled (\cref{tab:test:days}), and the sparse
day changes no ordering.

\paragraph{Results.}  \Cref{tab:test:main} gives the test numbers in the shape of
\cref{tab:horizons:synthesis}; \cref{tab:test:1h} gives every one-hour model of the study
with its validation and its test number side by side; the test days by day, the zero-shot
forecaster on the test samples and the value of the frozen offsets are in
\cref{app:results:test}.

\begin{table}[t]
\centering
\caption{Test days 24--26 March: the best model of each family at each horizon, all samples
and deduplicated samples (balanced accuracy), every model with its two logit offsets chosen
once on the validation day and applied unchanged.  ``Stateless'' is LightGBM refitted as in
validation and evaluated on the test samples of the stateful run; ``stateful'' is the best
run without trees evaluated on the test days (two passes with the engineered features at
1 min, 15 min, 1 h and 12 h; the time-first encoder at 6 h); ``residual'' is the two-pass
stateful model on the cross-fitted trees.  $\Delta$ in points against the stateless control,
95\% hour-block bootstrap intervals over the test hours.  TabM, not shown, is $+0.27$
[$+0.21$, $+0.34$] at 1 min, $+0.15$ [$+0.02$, $+0.27$] at 15 min and $-0.08$ [$-0.25$,
$+0.09$] at 1 h; a second seed of the one-minute stateful model is $+0.39$ [$+0.24$, $+0.57$]
on all samples and $+1.52$ [$+1.22$, $+1.83$] on the deduplicated ones, and a second seed
of the one-minute residual model $+0.96$ [$+0.85$, $+1.08$] and $+1.82$ [$+1.61$, $+2.02$].}
\label{tab:test:main}
\small
\begin{fittable}\begin{tabular}{llcccccc}
\toprule
Horizon & Subset & MR & Stateless & Stateful & $\Delta$ & Residual & $\Delta$ \\
\midrule
1 min  & All   & 0.5637 & 0.6496 & 0.6572 & $+0.76$ [$+0.60$, $+0.93$] & \textbf{0.6586} & $+0.90$ [$+0.78$, $+1.03$] \\
       & Dedup & 0.5342 & 0.6336 & 0.6443 & $+1.07$ [$+0.80$, $+1.34$] & \textbf{0.6539} & $+2.03$ [$+1.79$, $+2.26$] \\
15 min & All   & 0.5471 & 0.6088 & 0.6064 & $-0.25$ [$-0.44$, $-0.05$] & \textbf{0.6125} & $+0.36$ [$+0.24$, $+0.50$] \\
       & Dedup & 0.5301 & 0.6167 & 0.6090 & $-0.77$ [$-0.93$, $-0.60$] & \textbf{0.6234} & $+0.67$ [$+0.57$, $+0.77$] \\
1 h    & All   & 0.5268 & \textbf{0.5944} & 0.5806 & $-1.38$ [$-1.69$, $-1.07$] & 0.5933 & $-0.11$ [$-0.26$, $+0.03$] \\
       & Dedup & 0.5125 & 0.5887 & 0.5704 & $-1.83$ [$-1.96$, $-1.70$] & \textbf{0.5895} & $+0.08$ [$-0.02$, $+0.17$] \\
6 h    & All   & 0.5049 & \textbf{0.5851} & 0.5468 & $-3.83$ [$-4.11$, $-3.53$] & --- & \\
       & Dedup & 0.5025 & \textbf{0.5745} & 0.5323 & $-4.22$ [$-4.43$, $-4.00$] & --- & \\
12 h   & All   & 0.4724 & \textbf{0.5853} & 0.5539 & $-3.14$ [$-3.53$, $-2.73$] & 0.5780 & $-0.73$ [$-0.95$, $-0.50$] \\
       & Dedup & 0.4766 & \textbf{0.5749} & 0.5496 & $-2.53$ [$-2.85$, $-2.17$] & 0.5690 & $-0.59$ [$-0.75$, $-0.41$] \\
\bottomrule
\end{tabular}\end{fittable}
\end{table}

Three things carry over from the validation day to the test days, and one does not.

\emph{The ordering of the families, and the whole one-hour chain.}  From fifteen minutes on the
calibrated control is the best model or within its interval of the best, every learned
model stands where it stood on 13 March, and the sparse day moves no sign.  At one hour the
twelve runs of the study keep the order of \cref{ch:stateful,ch:graph,ch:external} to the
last step, and every gap is within a few tenths of its validation value: the two-pass
residual model and the tabular network tie the control, the two-pass model without trees is
a point and a half below it, the time-first encoder is a point and a half above the GRU
reference, the forecast features keep their sign on both seeds (three tenths on one, a few
hundredths on the other), the transaction graph is
below the reference, and the three seeds of the reference are within a quarter of a point of
one another.  At fifteen minutes the tabular network is slightly above the control and the
two-pass stateful model slightly below it, on the deduplicated samples by more, as in
validation, and the residual model is above the control by a third of a point on all samples
and two thirds on the deduplicated ones, with intervals that exclude zero, as in validation.  At one minute the gain of the sequential state is confirmed: the two-pass
stateful model is above the calibrated control on every subset and on each of the three
days, a second seed of the same model keeps the sign and the exclusion of zero on every
subset and every day while moving the size by four tenths of a point, which is the measure
of the seed noise at this horizon, and the residual model is above the control by about a
point on all samples and by two points on the deduplicated ones, with a second seed within
two tenths of the first.  The tabular network, the best stateless model at one minute, sits
between the control and the stateful models on every day: at one minute the stateful
models beat the strongest learner on the summary features, not only the trees.  At six and
twelve hours the LOB-only stateful models of the study are four to five points below the
control, their seeds within a point of one another in an order that changes on the sparse
day; at six hours the gap is half a point wider on the full stream and a point wider on the
deduplicated samples than in validation.  At twelve hours the two-pass model with the
engineered features is three points below the control and its residual on the trees within
a point of it, below with an interval that excludes zero on both views: the base closes most
of the gap on the test days as it did on the validation day.  The graph models stand where \cref{ch:graph} put them: below the LOB-only
reference at one hour, below the time-first model at six hours, and between the seeds of
the reference at twelve.

\begin{table}[t]
\centering
\caption{Every one-hour model of the study on the test days (balanced accuracy, all samples
and deduplicated samples, each model with its validation-day offsets; the validation column
is the validation day with the same offsets).  $\Delta$ in points against the calibrated
control, 95\% hour-block bootstrap intervals.}
\label{tab:test:1h}
\footnotesize
\begin{fittable}\begin{tabular}{lccccc}
\toprule
Model & Validation & Test & $\Delta$ [95\% CI] & Dedup. & $\Delta$ \\
\midrule
LightGBM (control)                           & 0.6044 & \textbf{0.5944} & --- & 0.5887 & --- \\
TabM, 3 seeds                                & 0.6036 & 0.5936 & $-0.08$ [$-0.25$, $+0.09$] & 0.5865 & $-0.22$ \\
MLP, 3 seeds                                 & 0.6002 & 0.5923 & $-0.22$ [$-0.40$, $-0.03$] & 0.5834 & $-0.54$ \\
Stateful, 2 passes, residual of the trees    & 0.6048 & 0.5933 & $-0.11$ [$-0.26$, $+0.03$] & \textbf{0.5895} & $+0.08$ \\
Stateful, 2 passes                           & 0.5907 & 0.5806 & $-1.38$ [$-1.69$, $-1.07$] & 0.5704 & $-1.83$ \\
Stateful, 1 pass, weight averaging           & 0.5879 & 0.5782 & $-1.62$ [$-1.97$, $-1.29$] & 0.5715 & $-1.72$ \\
Stateful, 1 pass                             & 0.5842 & 0.5738 & $-2.07$ [$-2.40$, $-1.73$] & 0.5667 & $-2.20$ \\
Time-first, no engineered features, seed 0         & 0.5718 & 0.5634 & $-3.11$ [$-3.48$, $-2.73$] & 0.5502 & $-3.85$ \\
Time-first, no engineered features, seed 1         & 0.5689 & 0.5604 & $-3.40$ [$-3.78$, $-3.03$] & 0.5483 & $-4.05$ \\
GRU reference + Chronos-2 features, seed 0   & 0.5598 & 0.5492 & $-4.52$ [$-4.92$, $-4.15$] & 0.5357 & $-5.30$ \\
GRU reference + Chronos-2 features, seed 1   & 0.5623 & 0.5511 & $-4.33$ [$-4.71$, $-3.95$] & 0.5350 & $-5.38$ \\
GRU reference, seed 0                        & 0.5569 & 0.5487 & $-4.57$ [$-5.00$, $-4.18$] & 0.5326 & $-5.61$ \\
GRU reference, seed 1                        & 0.5557 & 0.5480 & $-4.65$ [$-5.04$, $-4.26$] & 0.5314 & $-5.73$ \\
GRU reference, seed 2                        & 0.5561 & 0.5461 & $-4.83$ [$-5.23$, $-4.42$] & 0.5253 & $-6.34$ \\
GRU reference with the transaction graph     & 0.5543 & 0.5447 & $-4.97$ [$-5.39$, $-4.56$] & 0.5268 & $-6.19$ \\
Mean Reversion                               & 0.5311 & 0.5268 & $-6.76$ [$-7.23$, $-6.31$] & 0.5125 & $-7.62$ \\
\bottomrule
\end{tabular}\end{fittable}
\end{table}

\emph{The control.}  Its balanced accuracy is within a point and a half of its validation
value at every horizon, and at the long horizons it is higher on the test days than on the
validation day.

\emph{The rule.}  Mean Reversion keeps its ordering across the horizons, with the one-minute
reversal almost three points stronger on the test days and the twelve-hour one falling
below one half; the control's margin over the rule runs from six points at fifteen minutes
to eleven at twelve hours.  The zero-shot forecaster (\cref{tab:test:zeroshot}) behaves as on
the validation day: the sign of its median is far below the rule at every horizon, because it
almost never predicts \textsc{const}; with the half-tick dead zone it is one to three points
below the rule from one minute to one hour, ties it at six hours and is above it at twelve,
where the rule itself fell and the forecaster kept its validation value.  It re-discovers
the reversal, reads it a little more robustly than the rule on days when the twelve-hour
reversal weakens, and stays far below every learned model.

What does not carry over is the \emph{calibration} (\cref{tab:app:offsets}).  The offsets
chosen on 13 March, worth a few tenths of a point to the control there, cost it a few tenths
on the test days at fifteen minutes and one hour and are worth a few tenths at six and twelve
hours; at one minute, where they were worth more than five points in validation, they are
worth two.  For the stateful models the same frozen offsets help at every horizon, and most
where the model is worst calibrated: five to nine points at twelve hours, two and a half at
one minute.  A decision threshold fitted on one day is thus worth a few tenths of a point on
another day, in either direction, for a learner whose class mix is already right, and
several points for a model whose class mix is wrong; the comparisons of
\cref{tab:test:main,tab:test:1h} are between calibrated models and do not depend on it, but
any single raw number does, and the one-minute numbers most of all.

The test days therefore confirm the three regimes of \cref{sec:main:synthesis} on data that
no choice of the thesis had seen: at one minute the sequence beats the summary, at fifteen
minutes the stateful model on the trees is above the stateless learner, at one hour the
stateless learner, the tabular network and the stateful model on the trees tie, at six
hours the stateless learner is alone, and at twelve hours its residual is within a point of
it.
